\begin{ChineseAbstract}

建筑群能源系统同时受到天气、负荷、能源价格、分布式能源和用户舒适性需求的影响，具有强耦合、强不确定性和多目标特征。传统固定规则控制适应性不足，模型预测控制依赖精确模型并具有较高计算成本，而纯多智能体强化学习又容易面临样本效率低、训练不稳定和安全约束难以满足等问题。CityLearn为建筑群需求侧管理提供了统一的仿真环境、数据接口和评价指标，为控制策略的公平比较与可复现研究提供了基础。

在提出改进之前，现有方法仍面临三类核心难题：其一，以CHESCA为代表的规则与预测混合基线虽稳健可解释，但参数相对静态，难以及时适应负荷漂移、设备老化与极端天气；其二，纯多智能体强化学习虽具有长期优化潜力，却在小样本、非平稳环境和安全约束下难以稳定落地；其三，建筑群节能控制实验普遍存在难管理、难理解与难工程化使用的问题，算法配置、KPI对比与决策过程缺少统一可视化支撑。

针对上述难题，本文面向CityLearn建筑群能源管理场景，提出基于社区分层协调控制与残差多智能体强化学习的混合节能控制方法（CHESCA-ResMARL），并实现配套可视化管理平台。首先，系统梳理并复现CHESCA分层控制链路，建立标准化状态、动作与KPI模型；其次，在CHESCA基准动作之上叠加轻量级多智能体残差策略，通过动作掩码与安全投影实现“先验保安全、学习求增益”；再次，设计融合电价电费、制冷电费、电池套利、电池损耗与热舒适性的多目标奖励，并以消融、多种子与多场景实验评价方法有效性；最后，依托Vue、Spring Boot、MySQL与Python算法组件，贯通数据管理、算法配置、仿真任务、KPI分析与决策轨迹追踪。

项目已经完成CityLearn数据集构建、CHESCA基线评估、残差多智能体训练入口、动作安全约束、KPI导出和前后端可视化平台。实验结果表明，混合策略在保证舒适度的前提下实现了成本、碳排放、峰值与爬坡等多项指标的协同优化。本文的研究可为建筑群需求侧响应、低碳运行控制和智能能源管理系统的工程化应用提供参考。

\ChineseKeywords{建筑群节能；CityLearn；多智能体强化学习；残差强化学习；CHESCA；可视化管理}
\end{ChineseAbstract}

\begin{EnglishAbstract}

Building energy systems are affected by weather, demand, energy prices, distributed energy resources, and occupant comfort requirements. They therefore exhibit strong coupling, uncertainty, and multiple conflicting objectives. Fixed-rule control lacks adaptability, model predictive control relies on accurate models and high computational cost, while pure multi-agent reinforcement learning often suffers from low sample efficiency, training instability, and difficulty in satisfying safety constraints. CityLearn provides a standardized simulation environment, data interface, and evaluation metrics for building-community demand-side management, enabling reproducible and fair comparisons of control strategies.

Before the proposed improvements, existing approaches still face three core challenges. First, rule-and-prediction baselines such as CHESCA are robust and interpretable, yet relatively static and slow to adapt to load drift, equipment aging, and extreme weather. Second, pure multi-agent reinforcement learning has long-horizon optimization potential, but is difficult to deploy stably under limited samples, non-stationary environments, and hard safety constraints. Third, building-community energy-control experiments are commonly hard to manage, hard to interpret, and hard to engineer into usable systems, lacking unified support for algorithm configuration, KPI comparison, and decision tracing.

To address these challenges, this thesis investigates a hybrid energy-saving control method based on community hierarchical coordination and residual multi-agent reinforcement learning (CHESCA-ResMARL) for the CityLearn scenario, and develops an accompanying visualization management platform. First, the CHESCA hierarchical control pipeline is systematically reviewed and reproduced, and standardized state, action, and KPI models are established. Second, lightweight multi-agent residual policies are superimposed on CHESCA baseline actions, with action masks and safety projection realizing ``prior for safety, learning for gain''. Third, a multi-objective reward combining electricity cost, cooling cost, battery arbitrage, battery degradation, and thermal comfort is designed, and ablation, multi-seed, and cross-scenario evaluations are used to assess effectiveness. Finally, a visualization platform based on Vue, Spring Boot, MySQL, and Python algorithm services covers data management, algorithm configuration, simulation tasks, KPI analysis, and decision-trace inspection.

The current project has implemented CityLearn dataset construction, CHESCA baseline evaluation, a residual multi-agent training entry point, action safety constraints, KPI export, and a complete front-end/back-end visualization platform. Experimental results show that the hybrid strategy achieves coordinated optimization across cost, carbon emissions, peak load, and ramping while maintaining occupant comfort. This work may provide a reference for building-community demand response, low-carbon operation control, and engineering-oriented intelligent energy management systems.

\EnglishKeywords{building energy saving; CityLearn; multi-agent reinforcement learning; residual reinforcement learning; CHESCA; visualization management}
\end{EnglishAbstract}
